\documentclass[11pt]{article}
\usepackage[margin=1in]{geometry}
\usepackage{amsmath, amssymb}

\title{Dana Stewart Scott}
\author{}
\date{}

\begin{document}
\maketitle
\thispagestyle{empty}

Dana Stewart Scott (b.\ 1932) is one of the defining figures of twentieth-century
logic and theoretical computer science---a mathematician whose work reshaped how
we understand computation, meaning, and proof.

A student of \textbf{Alonzo Church} at Princeton (Ph.D., 1958), he first made his
mark in \textbf{mathematical logic and set theory}. His 1961 theorem that a
measurable cardinal contradicts G\"odel's axiom of constructibility---%
\emph{measurable cardinals imply $V \neq L$}---was a landmark, and
\emph{``Scott's trick''} for forming equivalence-class representatives without
the axiom of choice remains a standard tool. He developed Boolean-valued models
of set theory and neighborhood (``Scott--Montague'') semantics for modal logic.

In \textbf{1976 he received the Turing Award}, jointly with \textbf{Michael
Rabin}, for their early work on nondeterministic finite automata---a cornerstone
of the theory of computation.

He is perhaps best known for founding \textbf{domain theory} and, with
\textbf{Christopher Strachey}, the \textbf{denotational semantics} of programming
languages. His construction of $D_\infty$ gave the first mathematical model of
the untyped $\lambda$-calculus---a nontrivial space isomorphic to its own
function space, $D \cong [D \to D]$---and with it continuous lattices, Scott
domains, the Scott topology, and a rigorous meaning for recursion. It is fitting
that this project checks his own papers in Lean,\footnote{The ScottLean4 project
formalizes 38 of Dana Scott's papers (1958--2023) in Lean~4, machine-checked by
the kernel.} by the very standards of rigor he helped make possible.

Later honored with the \textbf{Rolf Schock Prize} in Logic and Philosophy (2001),
and long the Hillman University Professor at \textbf{Carnegie Mellon} (now
emeritus), he continues to work---recently on the mechanized formalization of
logic and category theory.

\end{document}
